% !PRINTED 263
\weberpage{263}

Betrachten wir nun von diesem Gesichtspunkte die Reihe der Grössen
\[
R_m,\ R_{m-1},\ R_{m-2} \ldots R_1,\ R_0,
\]
so findet beim Uebergang von $R_k$ zu $R_{k-1}$ eine Zeichenfolge oder ein Zeichenwechsel statt, je nachdem der Quotient $R_k : R_{k-1}$ positiv oder negativ ist.

Wir können also auch den folgenden Satz aussprechen:

V.\ \gesperrt{Ist $\pi$ die Anzahl der positiven, $\nu$ die der negativen Quadrate von $\varphi$, so ist $\pi$ gleich der Anzahl der Zeichenfolgen, $\nu$ gleich der Anzahl der Zeichenwechsel in der Kette}
\weq{4}{R_m,\ R_{m-1},\ R_{m-2} \ldots R_1,\ R_0.}

Diese Fassung des Satzes hat den Vorzug, dass sie sich auf den Fall übertragen lässt, dass in der Reihe (4) einzelne innere Glieder verschwinden, wenn nur nicht zwei auf einander folgende Glieder Null sind (was nach \S.\,82, III. immer angenommen werden kann). Wenn nämlich $R_k = 0$, $R_{k-1}$, $R_{k+1}$ von Null verschieden sind, so haben nach \S.\,82, Satz IV $R_{k-1}$ und $R_{k+1}$ verschiedene Vorzeichen. In der Reihe
\[
R_{k+1},\ R_k,\ R_{k-1}
\]
findet also ein Zeichenwechsel und eine Zeichenfolge statt, gleichviel ob wir das verschwindende $R_k$ durch eine positive oder eine negative Grösse ersetzen.

Wenn wir nun die Kette (4)
\[
R_m,\ R_{m-1} \ldots R_1,\ R_0,
\]
in der einzelne Glieder verschwinden, durch eine andere ersetzen,
\weq{5}{R'_m,\ R'_{m-1} \ldots R'_1,\ R'_0,}
in der kein Glied verschwindet, und in der den nicht verschwindenden Gliedern der Reihe (4) Glieder von demselben Vorzeichen entsprechen, so haben die Reihen (4) und (5) gleich viele Zeichenwechsel, welches Zeichen auch die den verschwindenden $R$ entsprechenden $R'$ haben mögen.

Es sei nun $\psi(x)$ eine zweite beliebige quadratische Form der Variablen $x$, mit der wir die Form
\weq{6}{\varphi' = \varphi + \varepsilon \psi}
bilden, worin $\varepsilon$ ein noch unbestimmter Co\"efficient ist. Wir werden nun sogleich zeigen, dass wir $\varepsilon$ so wählen können, dass $\varphi'$ und $\varphi$ dieselbe Zahl von positiven und negativen Quadraten
%CONT>>
